% ================================================
% 文件: 03_RFID技术.tex
% 章节: 第16章 RFID技术
% 所属部分: 系统与设备
% 版本: v1.0
% ================================================
% 本章目录结构（树状）:
% \chapter{RFID技术}
% ├── \section{RFID概述}
% ├── \section{RFID原理与分类}
% ├── \section{RFID系统组成}
% ├── \section{RFID应用场景}
% └── \section{RFID技术发展}
% ================================================

\chapter{RFID技术}
\label{chap:rfid_technology}

\section{RFID概述}
\label{sec:rfid_intro}

RFID（Radio Frequency Identificat
ion）即射频识别，是一种利用无线电波进行非接触式自动识别的技术
。RFID 能够识别目标对象并获取相关数据，无需人工干预，也无需
目标处于视线之内，因而在物流、零售、交通与制造等领域迅速普及。与
条形码相比，RFID 可同时读取多个标签、可穿过包装识别、并可写
入信息，适应性更强。

RFID 的主要特点可归纳如下：

\begin{itemize}
    \item \textbf{非接触识别}：无需物理接触或
对准即可读取数据，适合高速、批量场景。
    \item \textbf{快速识别}：配合防碰撞算法可
同时识别数十乃至上百个标签。
    \item \textbf{耐久性}：无源标签无电池、无
活动部件，可在潮湿、油污、灰尘等恶劣环境下长期工作。
    \item \textbf{大容量}：标签芯片可存储数十
字节至数 kB 数据，支持唯一标识与业务字段。
\end{itemize}

从系统价值看，RFID 将"物理对象"与"数字信息"绑定，是物联
网感知层的关键技术之一。它通过为每一物品赋予唯一电子身份，使采集
、追踪与决策自动化，显著降低人工成本并提升数据时效性。需要指出的
是，RFID 的识别距离与可靠性受频率、天线、环境介质（金属、液
体）与标签方向共同影响，工程部署需结合实际场景测算而非照搬标称值
。

与一维/二维条码相比，RFID 无需视线对准、可批量与穿透读取，
且标签可重写并嵌入安全机制；代价是单位成本较高、金属/液体环境影
响大，且需防范隐私追踪。RFID 的概念可追溯至二战敌我识别（I
FF），20 世纪 80 年代后随集成电路与天线技术进步走向商用
，21 世纪初在供应链管理（如零售 EPC 试点）中规模落地。今
天的 RFID 已与传感器、定位与通信深度融合，成为工业物联网中
不可或缺的自动识别手段之一。

\section{RFID原理与分类}
\label{sec:rfid_principle}

RFID 的工作原理是利用射频信号通过空间耦合实现无接触信息传递
与能量交换。按耦合机理可分为两类：\textbf{电感耦合}（近
场，磁耦合）与\textbf{电磁反向散射}（远场，电磁波反射）
，前者用于低频/高频近距系统，后者用于超高频/微波远距系统。

电感耦合依靠阅读器与标签天线间的互感。阅读器线圈产生交变磁场，标
签线圈感应出电动势为其芯片供电并解调指令；标签通过改变自身负载调
制（Load Modulation）反向改变磁场，将应答回传给阅
读器。其谐振频率由线圈电感 $L$ 与谐振电容 $C$ 决定：
\begin{equation}
f_0 = \frac{1}{2\pi\sqrt{LC}}.
\end{equation}
此类系统作用距离通常不超过线圈尺寸的数倍，典型为几厘米至一米。电
磁反向散射则利用电磁波传播：阅读器发射连续波，标签天线接收后由芯
片调制反射系数，将信息以"反向散射"形式回射；阅读距离近似服从雷
达距离方程的趋势，与发射功率、天线增益及波长相关，波长与频率关系
为
\begin{equation}
\lambda = \frac{c}{f},
\end{equation}
其中 $c\approx 3\times 10^8\ \math
rm{m/s}$ 为光速。频率越高波长越短、方向性越强、可获更远
识别距离，但对金属/液体的绕射与穿透能力下降。

按工作频率分类：

\begin{itemize}
    \item \textbf{低频（LF）}：$125\s
im 134\ \mathrm{kHz}$，识别距离短（通常 $
<1\ \mathrm{m}$），抗金属/液体好。
    \item \textbf{高频（HF）}：$13.56
\ \mathrm{MHz}$，识别距离中等（$1\sim 2\
 \mathrm{m}$），应用成熟（如门禁、图书）。
    \item \textbf{超高频（UHF）}：$860
\sim 960\ \mathrm{MHz}$，识别距离远（$3
\sim 10\ \mathrm{m}$），适合物流批量读取。
    \item \textbf{微波}：$2.45\ \ma
thrm{GHz}$、$5.8\ \mathrm{GHz}$，识
别距离最远，常用于 ETC、电子车牌。
\end{itemize}

按供电方式可分为有源 RFID（标签自带电池，距离远、寿命受限）
与无源 RFID（靠射频能量供电，免维护、寿命长），以及半有源（
电池仅供芯片、通信仍靠反向散射）。表\ref{tab:rfid_
band} 汇总了频率特性。

\begin{table}[htbp]
\centering
\caption{RFID 按工作频率分类对照}
\label{tab:rfid_band}
\begin{tabular}{|p{2.4cm}|p{3.0cm}|p{3.2cm}|p{4.0cm}|}
\hline
\textbf{频段} & \textbf{耦合方式} & \textbf{典型距离} & \textbf{典型应用} \\
\hline
LF $125\ \mathrm{kHz}$ & 电感耦合 & $<1\ \mathrm{m}$ & 动物标识、门禁 \\
\hline
HF $13.56\ \mathrm{MHz}$ & 电感耦合 & $1\sim 2\ \mathrm{m}$ & 图书、票证、支付 \\
\hline
UHF $860\sim 960\ \mathrm{MHz}$ & 反向散射 & $3\sim 10\ \mathrm{m}$ & 物流、仓储、零售 \\
\hline
微波 $2.45/5.8\ \mathrm{GHz}$ & 反向散射 & 最远 & ETC、电子车牌 \\
\hline
\end{tabular}
\end{table}

对于电感耦合系统，读取范围近似正比于线圈尺寸、品质因数与相对磁导
率，常用经验关系
\begin{equation}
d \approx k\sqrt{\frac{A_c Q}{f}},
\end{equation}
其中 $A_c$ 为线圈面积、$Q$ 为品质因数、$f$ 为工作
频率，$k$ 为与天线几何相关的常数。提高线圈 $Q$ 可增强耦
合但会收窄带宽，工程上需折衷。电磁反向散射系统属远场工作，读取距
离受发射功率、天线增益、标签雷达截面与环境多径共同影响，并要求阅
读器与标签极化匹配、方向对准；金属与液体会吸收或反射射频能量、改
变天线阻抗，导致读距骤降，常以特殊标签（如抗金属标签）或离物安装
缓解。

按耦合与频率选择体制时，还应权衡数据速率与抗冲突能力：LF/HF
 速率低但稳定、抗扰好，适合身份与支付；UHF/微波速率高、可群
读，适合物流与交通，但对环境更敏感。标签方向性在 UHF 尤其明
显——标签天线与阅读器极化垂直时会严重衰减，部署时常以多极化天线
或调整粘贴姿态来改善。

\section{RFID系统组成}
\label{sec:rfid_components}

一个完整的 RFID 系统由以下部分组成：

\begin{itemize}
    \item \textbf{标签（Tag）}：存储数据的
载体，由天线与芯片（ASIC）组成；含全球唯一序列号（UID）与
用户存储区。
    \item \textbf{阅读器（Reader）}：发
射激励、解调标签回波并解析数据，含射频前端、基带与接口。
    \item \textbf{天线}：发射与接收射频信号，
其增益、极化与波束影响识别区与方向性。
    \item \textbf{中间件}：对多阅读器数据进行
过滤、去重与事件封装，向上层提供统一接口。
    \item \textbf{应用系统}：数据库与业务软件
，完成数据分析、盘点与流程联动。
\end{itemize}

标签芯片是技术核心：无源芯片需以极低功耗完成整流取电、时钟恢复、
指令解析与负载调制，其灵敏度为系统读距的关键。阅读器则负责复杂的
防碰撞调度、密钥运算与协议栈处理。天线设计需兼顾阻抗匹配与近场均
匀性或远场覆盖，金属/液体环境常需特殊天线或隔离措施。中间件将原
始"读到/未读到"事件转化为业务语义（如入库、出库、越界），屏蔽
底层多厂商差异，是大规模部署的枢纽。阅读器接收标签反向散射信号的
功率可粗略表示为
\begin{equation}
P_r \propto \frac{P_t G_t G_r \lambda^2 \sigma}{(4\pi)^3 d^4},
\end{equation}
$\sigma$ 为标签雷达截面，$d$ 为距离；可见远场读距大
致随 $\sqrt[4]{P_t}$ 增长，提升发射功率对距离改
善有限而更易超标，故工程上更倚重天线与标签设计。

标签芯片的存储器通常划分为只读的 TID（标签唯一标识，含厂商与
型号）、可写的 EPC（电子产品代码）与 USER（用户数据）区
，并设有访问/摧毁口令保护。无源芯片以整流电路从射频场取电，其"
启动灵敏度"（唤醒所需最小场强）直接决定读距；为降低功耗，芯片多
采用时钟门控与低电压工艺。阅读器则由射频前端（发射链、接收链、环
行器/耦合器）、基带处理器与网络接口构成，承担编码调制、防碰撞调
度、解密与协议栈处理，其接收灵敏度（能解调标签微弱反向散射的最小
信号）是远场系统性能瓶颈之一。中间件对多阅读器上报做过滤、去重与
平滑，屏蔽重复读取与抖动，向上提供"出现/消失/移动"等业务事件
。

\section{RFID应用场景}
\label{sec:rfid_applications}

RFID 的应用已从识别拓展到全流程感知：

\begin{itemize}
    \item \textbf{物流管理}：仓储管理、货物追
踪、供应链管理，实现出入库自动核对与可视化。
    \item \textbf{零售行业}：库存盘点、防盗（
EAS）、自助结账，降低缺货与损耗。
    \item \textbf{交通领域}：ETC 不停车收
费、车辆识别、停车场管理，提升通行效率。
    \item \textbf{医疗领域}：药品溯源、患者识
别、医疗器械追踪，减少用药与标本差错。
    \item \textbf{身份识别}：门禁、考勤、电子
证照，兼顾便捷与安全。
\end{itemize}

典型应用剖析——以超高频仓储盘点为例：在仓库出入口与货位布置多天
线阅读器，叉车或通道批量读取托盘/箱体标签，中间件实时比对计划单
生成差异告警，实现"秒级"全盘与错发拦截。又如 ETC：车载 O
BU（微波标签）进入车道微波场时被路侧单元（RSU）唤醒，双方完
成双向认证与扣费握手，车辆无需停车即可通行。这些案例的共同点是把
"人工扫码"升级为"环境自动感知"，从而把数据时效性从分钟级压缩
到秒级以内，并减少人为错漏。

在物流场景中，RFID 取代人工扫码实现整托盘/整箱级自动采集，
配合传送带与门禁式阅读器完成出入库核验，显著降低错漏与人力；零售
门店借 RFID 实现"即拿即走"结算与实时库存，减少缺货与盘点
成本；交通领域除 ETC 外，RFID 还用于集装箱与货车电子关
锁、停车场不停车出入；医疗上用于标本/器械/药品全生命周期追踪，
防止用错与流失；身份识别则覆盖门禁、考勤、电子证照与赛事计时。这
些应用的价值不单是"识别"，更在于把物理对象的状态变化实时映射为
可分析的数据流，支撑精益管理与合规追溯。

\section{RFID技术发展}
\label{sec:rfid_development}

RFID 技术的发展趋势体现在性能、成本、安全与融合多个维度：

\begin{itemize}
    \item \textbf{更高频率与更远距离}：向 U
HF/微波与多天线波束成形发展，提升批量读取速度与覆盖。
    \item \textbf{更小尺寸与更低成本}：标签向
印刷电子、柔性衬底演进，单位成本持续下降以支撑海量粘贴。
    \item \textbf{更低功耗}：芯片电路优化与能
量收集，延长半有源标签寿命。
    \item \textbf{更好安全性}：引入轻量级加密
、双向认证与不可逆销毁，缓解伪造与追踪。
    \item \textbf{物联网集成}：与传感、定位、
边缘计算融合，从"识别"走向"感知+决策"。
\end{itemize}

在标准化方面，EPCglobal 提出的电子产品代码（EPC）体
系为物品定义了全球唯一的编码结构与信息网络标识，配合 EPCIS
 事件标准实现跨企业共享；ISO/IEC 18000 系列则按频
段规范了空中接口（如 ISO/IEC 18000-3 对应 HF
、18000-6 对应 UHF），为不同厂商设备互通奠定基础。防
碰撞是多标签读取的核心算法：基于概率的 ALOHA 类（含时隙 
ALOHA、帧时隙 ALOHA）实现简单、适于动态标签群；基于确
定性的二进制树（Binary Tree）类（含 Query 树、
动态二进制树）可保证在无碰撞时完全清点、吞吐稳定。随标签密度上升
，自适应混合算法成为主流。

安全与隐私是 RFID 规模化不可回避的问题：无源标签算力有限，
难以承载重型公钥运算，故多采用轻量级对称认证与单向哈希；对"位置
隐私"可通过静默/销毁机制或每次会话刷新假名（Pseudo-ID
）缓解。总之，RFID 正从孤立的识别手段，演化为可感知、可认证
、可联网的泛在标识基础设施，并与条码、传感器、蜂窝与卫星链路互补
，共同支撑数字孪生与全程可追溯的新型产业体系。

防碰撞是多标签并发读取的核心算法。基于随机退避的帧时隙 ALOH
A 将时间分为若干时隙，每标签随机选一时隙应答，当标签数 $n$
 与时隙数 $N$ 匹配时吞吐接近最优；其理论吞吐约为
\begin{equation}
S = \frac{n}{N}\left(1-\frac{1}{N}\right)^{n-1},
\end{equation}
在 $n\approx N$ 时取得峰值约 $1/e\appro
x 36.8\%$。确定性二进制树算法（如 Query 树、动态
二进制树）通过逐位分裂冲突前缀保证在有限轮次内完整清点、无漏读，
吞吐随标签数更平稳但交互开销较大；实际系统常按标签密度在两类间自
适应切换。标准化方面，EPCglobal 定义了 EPC 编码、
标签级信息（TDS）与事件共享（EPCIS），ISO/IEC 1
8000 系列按频段规定空中接口（18000-2 为 LF、18
000-3 为 HF、18000-6 为 UHF、18000-7
 为 $433\ \mathrm{MHz}$、18000-4 为
 $2.45\ \mathrm{GHz}$），ISO/IEC 1
4443 与 15693 分别规范近距与疏耦合卡，为跨厂商互通与
合规检测奠定基础。未来 RFID 将与印刷电子、能量收集、传感融
合及边缘智能结合，从"给物发身份证"走向"让物自己说话"，在智能
制造、无人零售与城市治理中发挥更大作用。
